\chapter[Camino descartado: Contact ID]{Camino evaluado y descartado: el protocolo Contact ID}
\label{anx:contactid}
\small

Durante la fase inicial de esta etapa (17 de agosto al 1 de septiembre) el trabajo evaluó como eje el
protocolo \textbf{Contact ID (SIA DC-05)} de la industria de alarmas, sobre transporte DC-09. Se
verificó empíricamente su gramática y su suma de comprobación con una implementación propia, se
evaluaron ocho alternativas mediante una matriz ponderada y se diseñó un perfil de aplicación. Ese
material completo permanece en el repositorio del proyecto (\texttt{docs/05-protocolo-contactid.md} y
\texttt{plataforma/herramientas/contactid.py}).

\textbf{Razón del descarte (reunión con el profesor guía, 2 de septiembre de 2026):} Contact ID está
orientado a \textbf{eventos discretos} —alarmas que ocurren o no— mientras que el problema de este
trabajo es \textbf{telemetría continua}: lecturas periódicas cada 5 a 30 minutos cuyo valor está en
la serie, no en el evento. El mensaje de 16 dígitos no transporta valores medidos de forma nativa, y
las extensiones necesarias para hacerlo desplazaban al protocolo fuera de su dominio de diseño.

\textbf{Lo que el análisis aportó y se conserva:} la restricción de que el silencio de un sensor
debe detectarse activamente (los códigos de supervisión de la industria de alarmas inspiran los
requerimientos RF-24 a RF-30, hoy resueltos con el \emph{keep-alive} y el testamento de MQTT); la
disciplina de evaluar alternativas con criterios ponderados y de verificar afirmaciones con código; y
la separación lógica entre eventos y telemetría, que sobrevive dentro del modelo interno.

El registro de la decisión, con las tareas derivadas, está en
\texttt{bitacora/2026-09-02-reunion-pivote.md}.
